\documentclass[12pt]{article}

\usepackage[utf8]{inputenc}
\usepackage{amsmath, amssymb, amsthm}
\usepackage{geometry}
\geometry{a4paper, margin=2.5cm}

\title{Probabilitate}
\author{}
\date{}

\theoremstyle{definition}
\newtheorem{definition}{Definiție}

\begin{document}

\maketitle

\begin{definition}
Fie $(\Omega, \mathcal{F})$ un spațiu măsurabil. 
O funcție $\mathbb{P} : \mathcal{F} \to [0,1]$ se numește \emph{probabilitate} dacă:
\begin{enumerate}
    \item $\mathbb{P}(\Omega) = 1$;
    \item pentru orice familie de mulțimi disjuncte două câte două $(A_n)_{n \ge 1}$ din $\mathcal{F}$ avem
    

\[
        \mathbb{P}\!\left( \bigcup_{n=1}^{\infty} A_n \right)
        = \sum_{n=1}^{\infty} \mathbb{P}(A_n).
    \]


\end{enumerate}
\end{definition}

\noindent
Tripletul $(\Omega, \mathcal{F}, \mathbb{P})$ se numește \emph{spațiu de probabilitate}.

\end{document}
